\documentclass[12pt]{article}
\usepackage[ukrainian]{babel}
\usepackage{fontspec}
% --- НАЛАШТУВАННЯ ШРИФТІВ ---
\setmainfont{Schoolbook-Regular.otf}[
    Path = ../../,
    BoldFont = Schoolbook-Bold.otf,
    ItalicFont = Schoolbook-Italic.otf,
    BoldItalicFont = Schoolbook-BoldItalic.otf
]
\usepackage[italic]{mathastext}
\usepackage[a4paper,margin=1.8cm,bottom=2cm,top=2cm]{geometry}
\usepackage{amsmath,amssymb}
\usepackage{tikz}
\usepackage{pgfplots}
\pgfplotsset{compat=1.16}
\usetikzlibrary{calc,patterns,angles,quotes}
\usepackage{xcolor,array,fancyhdr,enumitem,multicol}
\usepackage{colortbl}
\usepackage{tcolorbox}
\tcbuselibrary{skins, breakable}
\usepackage[hidelinks]{hyperref}
\usepackage[firstpage=false, text=@pvtr2525, scale=0.6, color=gray!12]{draftwatermark}
\definecolor{mainGreen}{RGB}{34, 120, 64}
\definecolor{yearOrange}{RGB}{220, 100, 30}
\definecolor{yearcolor}{RGB}{220, 100, 30}
\definecolor{titleBg}{RGB}{225, 240, 220}
\definecolor{instrBg}{RGB}{255, 240, 220}
\definecolor{instrBorder}{RGB}{220, 150, 80}
\pagestyle{fancy}
\fancyhf{}
\renewcommand{\headrulewidth}{0.4pt}
\renewcommand{\footrulewidth}{0.4pt}
\fancyhead[C]{\small\color{gray!80} Збірник НМТ \textendash{} сесії 23.05--5.06.2026}
\fancyhead[R]{\small\color{mainGreen}\textbf{\href{https://t.me/pvtr2525}{Pan Markiyan}}}
\fancyfoot[L]{\small\color{gray!80}\href{https://t.me/pvtr2525}{ТГ: @pvtr2525}}
\fancyfoot[C]{\small\color{gray!80}\thepage}
\fancyfoot[R]{\small\color{gray!80}НМТ 2026}
\newcommand{\answerTable}[5]{%
\par\vspace{0.15cm}
\noindent
\begin{tabular}{|*{5}{>{\centering\arraybackslash}m{3cm}|}}
\hline
\rule[-0.15cm]{0pt}{0.6cm}\textbf{А} & \rule[-0.15cm]{0pt}{0.6cm}\textbf{Б} & \rule[-0.15cm]{0pt}{0.6cm}\textbf{В} & \rule[-0.15cm]{0pt}{0.6cm}\textbf{Г} & \rule[-0.15cm]{0pt}{0.6cm}\textbf{Д} \\
\hline
\rule[-0.2cm]{0pt}{0.75cm}#1 & \rule[-0.2cm]{0pt}{0.75cm}#2 & \rule[-0.2cm]{0pt}{0.75cm}#3 & \rule[-0.2cm]{0pt}{0.75cm}#4 & \rule[-0.2cm]{0pt}{0.75cm}#5 \\
\hline
\end{tabular}
\par\vspace{0.25cm}
}
\newcommand{\sectionTitle}[1]{%
\par\vspace{0.4cm}
\begin{tcolorbox}[colback=titleBg, colframe=titleBg, boxrule=0pt, arc=8pt,
    halign=center, fontupper=\Large\bfseries\color{mainGreen}]
#1
\end{tcolorbox}
\par\vspace{0.3cm}
}
\newcommand{\nmtAnswerBox}{%
\par\vspace{0.2cm}\noindent
Відповідь:\ \
\begin{tabular}{@{}*{4}{|>{\centering\arraybackslash}p{0.8cm}}|@{\hspace{0.1cm}\textbf{,}\hspace{0.1cm}}*{3}{|>{\centering\arraybackslash}p{0.8cm}}|@{}}
\hline
\rule[-0.2cm]{0pt}{0.7cm} & & & & & & \\
\hline
\end{tabular}
\par\vspace{0.3cm}
}
\newcommand{\taskBlock}[1]{%
\par\vspace{0.45cm}
\noindent{\color{mainGreen}\rule{\linewidth}{0.8pt}}\par\vspace{0.12cm}
\noindent{\small\bfseries\color{mainGreen}#1}\par\vspace{0.25cm}
}
\newcounter{zad}
\newcommand{\zadtask}[1]{\stepcounter{zad}\noindent\textbf{\thezad.}\ \ #1\par\vspace{0.2cm}}
\newcommand{\zadnum}{\stepcounter{zad}\textbf{\thezad.}\ \ }
\newcounter{ans}
\newcommand{\ansitem}[1]{\stepcounter{ans}\noindent\textbf{\theans.}\ #1\par\vspace{0.07cm}}
\newcommand{\nmtyear}[1]{\hfill{\small\color{yearOrange}(НМТ #1)}}
\setcounter{tocdepth}{2}
\newcommand{\chapterTitle}[1]{%
\clearpage\phantomsection\addcontentsline{toc}{section}{#1}%
\sectionTitle{#1}}
\newcommand{\typeTitle}[1]{%
\phantomsection\addcontentsline{toc}{subsection}{\quad #1}%
\par\vspace{0.3cm}
\noindent{\large\bfseries\color{yearOrange}#1}\par\vspace{0.08cm}
\noindent{\color{yearOrange}\rule{\linewidth}{0.4pt}}\par\vspace{0.2cm}}
\newcommand{\subTypeTitle}[1]{%
\phantomsection\addcontentsline{toc}{subsubsection}{\qquad #1}%
\par\vspace{0.25cm}
\noindent{\normalsize\bfseries\color{mainGreen!80!black}\textit{#1}}\par\vspace{0.05cm}
\noindent{\color{mainGreen!80!black}\rule{0.7\linewidth}{0.3pt}}\par\vspace{0.15cm}}
\newcommand{\ansTheme}[1]{\par\vspace{0.25cm}\noindent{\bfseries\color{mainGreen}#1}\par\vspace{0.12cm}}
\newcommand{\ansType}[1]{\par\vspace{0.1cm}\noindent{\itshape\color{yearOrange}#1}\par\vspace{0.08cm}}
% --- КНОПКА РОЗВ'ЯЗКУ ---
\newcommand{\solbtn}[1]{\hypertarget{task:#1}{}\par\vspace{0.05cm}\noindent\hyperlink{sol:#1}{\fbox{\footnotesize розв'язок $\rightarrow$}}\par\vspace{0.2cm}}
\begin{document}
\catcode`\"=12\relax
\thispagestyle{empty}
\vspace*{1.2cm}
\begin{center}
{\Large\bfseries\color{mainGreen} ЗБІРНИК НМТ-2026 з математики}\\[0.4cm]
{\Huge\bfseries\color{mainGreen} РОЗДІЛ 5. СТЕРЕОМЕТРІЯ}\\[0.6cm]
{\large Оригінали сесій 23.05--5.06.2026 з короткими розв'язками та авторськими аналогами}\\[1cm]
{\large\color{mainGreen}\textbf{\href{https://t.me/pvtr2525}{Телеграм: @pvtr2525}}}
\end{center}
\clearpage
\phantomsection\addcontentsline{toc}{section}{Зміст}
\sectionTitle{ЗМІСТ}
{\hypersetup{linkcolor=black}\renewcommand{\contentsname}{}\vspace{-1cm}\tableofcontents}
\chapterTitle{РОЗДІЛ 5. СТЕРЕОМЕТРІЯ}
\noindent У цьому розділі~--- $66$ задач зі стереометрії з усіх $7$ сесій НМТ-2026 (23.05--5.06). Задачі згруповано у три підтеми:
\par\vspace{0.25cm}
\noindent\textbullet\ \textit{Просторові фігури та їх властивості}~--- теоретичні питання про перетин площин, паралельність і мимобіжність прямих, апофеми й ребра пірамід, тіла обертання, дотичні до сфери.
\par\vspace{0.15cm}
\noindent\textbullet\ \textit{Вектори і координати в просторі}~--- модулі векторів, координати точок і середин, відстані, рівняння сфери, координати в задачах про циліндр.
\par\vspace{0.15cm}
\noindent\textbullet\ \textit{Стереометричні розрахунки (коротка відповідь)}~--- комбінації двох тіл (циліндр+конус, призма+конус, призма+піраміда, циліндр+паралелепіпед, конус+півсфера) з обчисленням об'єму чи поверхні.
\par\vspace{0.25cm}
\noindent Під кожною оригінальною задачею розміщено кнопку \fbox{\footnotesize розв'язок $\rightarrow$}, що веде до коментаря в кінці розділу.
\par
\typeTitle{Завдання з вибором однієї відповіді (А--Д)}
\subTypeTitle{Просторові фігури та їх властивості}
\taskBlock{Перетин площин у піраміді --- сесія 23.05, завдання №3}
\noindent
\begin{minipage}[c]{0.58\textwidth}
\zadnum У чотирикутній піраміді $SABCD$ точка $M$ належить ребру $AB$. По якій прямій перетинаються площини $ABS$ і $MCS$? \nmtyear{2026}
\end{minipage}%
\hfill
\begin{minipage}[c]{0.40\textwidth}
\centering
\begin{tikzpicture}[scale=0.85, thick, line join=round]
    \coordinate (A) at (1.6,0);
    \coordinate (B) at (-0.8,1.0);
    \coordinate (C) at (1.6,1.0);
    \coordinate (D) at (3.4,0.3);
    \coordinate (S) at (1.1,3.2);
    \coordinate (M) at ($(A)!0.55!(B)$);
    \draw (A) -- (S) -- (D) -- cycle;
    \draw (A) -- (D);
    \draw (S) -- (C);
    \draw (C) -- (D);
    \draw[dashed] (B) -- (C);
    \draw (B) -- (S);
    \draw (A) -- (B);
    \fill (M) circle (1.5pt);
    \node[above] at (S) {$S$};
    \node[left] at (B) {$B$};
    \node[above right] at (C) {$C$};
    \node[below] at (A) {$A$};
    \node[right] at (D) {$D$};
    \node[below left] at (M) {$M$};
\end{tikzpicture}
\end{minipage}
\par\vspace{0.2cm}
\answerTable{$MC$}{$AB$}{$SB$}{$BC$}{$SM$}
\solbtn{1}
\vspace{0.3cm}
\noindent
\begin{minipage}[c]{0.58\textwidth}
\zadnum У трикутній піраміді $SABC$ точка $K$ належить ребру $AC$. По якій прямій перетинаються площини $ACS$ і $KBS$?
\end{minipage}%
\hfill
\begin{minipage}[c]{0.40\textwidth}
\centering
\begin{tikzpicture}[scale=0.85, thick, line join=round]
    \coordinate (A) at (0,0);
    \coordinate (C) at (3.4,0.4);
    \coordinate (B) at (2.2,1.4);
    \coordinate (S) at (1.6,3.2);
    \coordinate (K) at ($(A)!0.5!(C)$);
    \draw (A) -- (C) -- (S) -- cycle;
    \draw (C) -- (B) -- (S);
    \draw[dashed] (A) -- (B);
    \draw[dashed] (B) -- (K);
    \fill (K) circle (1.5pt);
    \node[above] at (S) {$S$};
    \node[below left] at (A) {$A$};
    \node[right] at (B) {$B$};
    \node[below right] at (C) {$C$};
    \node[below] at (K) {$K$};
\end{tikzpicture}
\end{minipage}
\par\vspace{0.2cm}
\answerTable{$KB$}{$AC$}{$SK$}{$SB$}{$SA$}
\vspace{0.3cm}
\noindent
\begin{minipage}[c]{0.58\textwidth}
\zadnum У чотирикутній піраміді $SABCD$ точка $P$ належить ребру $CD$. По якій прямій перетинаються площини $SCD$ і $PBS$?
\end{minipage}%
\hfill
\begin{minipage}[c]{0.40\textwidth}
\centering
\begin{tikzpicture}[scale=1, thick, line join=round]
    \coordinate (A) at (1.6,0);
    \coordinate (B) at (-0.8,1.0);
    \coordinate (C) at (1.6,1.0);
    \coordinate (D) at (3.4,0.3);
    \coordinate (S) at (1.1,3.2);
    \coordinate (P) at ($(C)!0.5!(D)$);
    \draw (A) -- (S) -- (D) -- cycle;
    \draw (A) -- (D);
    \draw (S) -- (C);
    \draw (C) -- (D);
    \draw[dashed] (B) -- (C);
    \draw (B) -- (S);
    \draw (A) -- (B);
    \draw[dashed] (B) -- (P);
    \fill (P) circle (1.5pt);
    \node[above] at (S) {$S$};
    \node[left] at (B) {$B$};
    \node[above] at (C) {$C$};
    \node[below] at (A) {$A$};
    \node[right] at (D) {$D$};
    \node[above right] at (P) {$P$};
\end{tikzpicture}
\end{minipage}
\par\vspace{0.2cm}
\answerTable{$BP$}{$CD$}{$SC$}{$SP$}{$SD$}
\taskBlock{Апофеми піраміди --- сесія 30.05, завдання №2}
\noindent
\begin{minipage}[c]{0.60\textwidth}
\zadnum На рисунку зображено правильну трикутну піраміду $SABC$ з вершиною $S$. Скільки всього апофем можна провести в цій піраміді, якщо у неї $4$ грані? \nmtyear{2026}
\end{minipage}%
\hfill
\begin{minipage}[c]{0.38\textwidth}
\centering
\begin{tikzpicture}[scale=0.9, thick, line join=round]
    \coordinate (S) at (1.2, 3);
    \coordinate (A) at (1.5, 0);
    \coordinate (B) at (3, 1);
    \coordinate (C) at (0, 1);
    \draw[thick] (A) -- (B) -- (S) -- cycle;
    \draw[thick] (S) -- (A);
    \draw[thick, dashed] (C) -- (B);
    \draw[thick] (C) -- (A);
    \draw[thick] (C) -- (S);
    \node[above] at (S) {$S$};
    \node[below] at (A) {$A$};
    \node[right] at (B) {$B$};
    \node[left] at (C) {$C$};
\end{tikzpicture}
\end{minipage}
\par\vspace{0.2cm}
\begin{flushleft}
\textbf{А} \quad жодної \\[0.1cm]
\textbf{Б} \quad лише одну \\[0.1cm]
\textbf{В} \quad лише дві \\[0.1cm]
\textbf{Г} \quad лише три \\[0.1cm]
\textbf{Д} \quad більше трьох
\end{flushleft}
\solbtn{4}
\vspace{0.3cm}
\noindent
\begin{minipage}[c]{0.60\textwidth}
\zadnum На рисунку зображено правильну чотирикутну піраміду $SABCD$ з вершиною $S$. Скільки всього апофем можна провести в цій піраміді, якщо у неї $5$ граней?
\end{minipage}%
\hfill
\begin{minipage}[c]{0.38\textwidth}
\centering
\begin{tikzpicture}[scale=0.9, thick, line join=round]
    \coordinate (S) at (1.6, 3.2);
    \coordinate (A) at (0.2, 0);
    \coordinate (B) at (3.0, 0);
    \coordinate (C) at (3.6, 1.2);
    \coordinate (D) at (0.8, 1.2);
    \draw[thick] (A) -- (B) -- (S) -- cycle;
    \draw[thick] (S) -- (A);
    \draw[thick] (B) -- (C) -- (S);
    \draw[thick, dashed] (C) -- (D);
    \draw[thick, dashed] (D) -- (A);
    \draw[thick, dashed] (D) -- (S);
    \node[above] at (S) {$S$};
    \node[below left] at (A) {$A$};
    \node[below right] at (B) {$B$};
    \node[right] at (C) {$C$};
    \node[left] at (D) {$D$};
\end{tikzpicture}
\end{minipage}
\par\vspace{0.2cm}
\begin{flushleft}
\textbf{А} \quad жодної \\[0.1cm]
\textbf{Б} \quad лише дві \\[0.1cm]
\textbf{В} \quad лише три \\[0.1cm]
\textbf{Г} \quad лише чотири \\[0.1cm]
\textbf{Д} \quad більше чотирьох
\end{flushleft}
\vspace{0.3cm}
\noindent
\begin{minipage}[c]{0.60\textwidth}
\zadnum На рисунку зображено правильну шестикутну піраміду $SABCDEF$ з вершиною $S$. Скільки всього ребер має ця піраміда, якщо у неї $7$ граней?
\end{minipage}%
\hfill
\begin{minipage}[c]{0.38\textwidth}
\centering
\begin{tikzpicture}[scale=0.85, thick, line join=round]
    \coordinate (S) at (1.8, 3.4);
    \coordinate (A) at (0.0, 0.4);
    \coordinate (B) at (1.2, 0.0);
    \coordinate (C) at (3.0, 0.2);
    \coordinate (D) at (3.6, 0.9);
    \coordinate (E) at (2.4, 1.3);
    \coordinate (F) at (0.6, 1.1);
    \draw[thick] (A) -- (B) -- (C) -- (D);
    \draw[thick, dashed] (D) -- (E) -- (F) -- (A);
    \draw[thick] (S) -- (A);
    \draw[thick] (S) -- (B);
    \draw[thick] (S) -- (C);
    \draw[thick] (S) -- (D);
    \draw[thick, dashed] (S) -- (E);
    \draw[thick, dashed] (S) -- (F);
    \node[above] at (S) {$S$};
    \node[left] at (A) {$A$};
    \node[below] at (B) {$B$};
    \node[below right] at (C) {$C$};
    \node[right] at (D) {$D$};
    \node[above right] at (E) {$E$};
    \node[above left] at (F) {$F$};
\end{tikzpicture}
\end{minipage}
\par\vspace{0.2cm}
\begin{flushleft}
\textbf{А} \quad лише три \\[0.1cm]
\textbf{Б} \quad лише чотири \\[0.1cm]
\textbf{В} \quad лише п'ять \\[0.1cm]
\textbf{Г} \quad лише шість \\[0.1cm]
\textbf{Д} \quad більше шести
\end{flushleft}
\vspace{0.3cm}
\taskBlock{Мимобіжні прямі у призмі --- сесія 1.06, завдання №2}
\noindent
\begin{minipage}[c]{0.58\textwidth}
\zadnum У прямокутній чотирикутній призмі $ABCDA_1B_1C_1D_1$ проведено діагональ $B_1D$. Скільки ребер цієї призми \textit{мимобіжні} з прямою $B_1D$? \nmtyear{2026}
\end{minipage}%
\hfill
\begin{minipage}[c]{0.40\textwidth}
\centering
\begin{tikzpicture}[scale=0.6, thick, line join=round]
    \coordinate (A) at (0,0);
    \coordinate (D) at (3,0);
    \coordinate (B) at (1,1.5);
    \coordinate (C) at (4,1.5);
    \coordinate (A1) at (0,4);
    \coordinate (D1) at (3,4);
    \coordinate (B1) at (1,5.5);
    \coordinate (C1) at (4,5.5);
    \draw[dashed] (A) -- (B);
    \draw[dashed] (B) -- (C);
    \draw[dashed] (B) -- (B1);
    \draw[thick] (A) -- (D) -- (C) -- (C1) -- (D1) -- (A1) -- cycle;
    \draw[thick] (A1) -- (B1) -- (C1);
    \draw[thick] (D) -- (D1);
    \draw[thick] (A) -- (A1);
    \draw[mainGreen, line width=1.2pt] (B1) -- (D);
    \node[left] at (A) {$A$};
    \node[right] at (D) {$D$};
    \node[left] at (B) {$B$};
    \node[right] at (C) {$C$};
    \node[left] at (A1) {$A_1$};
    \node[right] at (D1) {$D_1$};
    \node[left] at (B1) {$B_1$};
    \node[right] at (C1) {$C_1$};
\end{tikzpicture}
\end{minipage}
\par\vspace{0.2cm}
\answerTable{$4$}{$5$}{$6$}{$7$}{$8$}
\solbtn{7}
\noindent
\begin{minipage}[c]{0.58\textwidth}
\zadnum У правильній трикутній призмі $ABCA_1B_1C_1$ проведено діагональ бічної грані $AB_1$. Скільки ребер цієї призми \textit{мимобіжні} з прямою $AB_1$?
\end{minipage}%
\hfill
\begin{minipage}[c]{0.40\textwidth}
\centering
\begin{tikzpicture}[scale=0.6, thick, line join=round]
    \coordinate (A) at (0,0);
    \coordinate (B) at (3.4,0);
    \coordinate (C) at (2.2,1.4);
    \coordinate (A1) at (0,4);
    \coordinate (B1) at (3.4,4);
    \coordinate (C1) at (2.2,5.4);
    \draw[dashed] (A) -- (C);
    \draw[dashed] (B) -- (C);
    \draw[dashed] (C) -- (C1);
    \draw[thick] (A) -- (B) -- (B1) -- (A1) -- cycle;
    \draw[thick] (A1) -- (C1) -- (B1);
    \draw[mainGreen, line width=1.2pt] (A) -- (B1);
    \node[below left] at (A) {$A$};
    \node[below right] at (B) {$B$};
    \node[right] at (C) {$C$};
    \node[above left] at (A1) {$A_1$};
    \node[above right] at (B1) {$B_1$};
    \node[above] at (C1) {$C_1$};
\end{tikzpicture}
\end{minipage}
\par\vspace{0.2cm}
\answerTable{$2$}{$3$}{$4$}{$5$}{$6$}
\noindent
\begin{minipage}[c]{0.58\textwidth}
\zadnum У кубі $ABCDA_1B_1C_1D_1$ проведено діагональ основи $AC$. Скільки ребер цього куба \textit{мимобіжні} з прямою $AC$?
\end{minipage}%
\hfill
\begin{minipage}[c]{0.40\textwidth}
\centering
\begin{tikzpicture}[scale=0.6, thick, line join=round]
    \coordinate (A) at (0,0);
    \coordinate (B) at (3,0);
    \coordinate (C) at (4,1.5);
    \coordinate (D) at (1,1.5);
    \coordinate (A1) at (0,3);
    \coordinate (B1) at (3,3);
    \coordinate (C1) at (4,4.5);
    \coordinate (D1) at (1,4.5);
    \draw[dashed] (A) -- (D);
    \draw[dashed] (D) -- (C);
    \draw[dashed] (D) -- (D1);
    \draw[thick] (A) -- (B) -- (C) -- (C1) -- (B1) -- (A1) -- cycle;
    \draw[thick] (A1) -- (D1) -- (C1);
    \draw[thick] (B) -- (B1);
    \draw[mainGreen, line width=1.2pt, dash pattern=on 4pt off 2pt] (A) -- (C);
    \node[below left] at (A) {$A$};
    \node[below right] at (B) {$B$};
    \node[right] at (C) {$C$};
    \node[left] at (D) {$D$};
    \node[left] at (A1) {$A_1$};
    \node[right] at (B1) {$B_1$};
    \node[above right] at (C1) {$C_1$};
    \node[above left] at (D1) {$D_1$};
\end{tikzpicture}
\end{minipage}
\par\vspace{0.2cm}
\answerTable{$4$}{$5$}{$6$}{$7$}{$8$}
\taskBlock{Паралельні прямі в кубі --- сесія 2.06, завдання №5}
\noindent
\begin{minipage}[c]{0.56\textwidth}
\zadnum На рисунку зображено куб $ABCDA_1B_1C_1D_1$. Скільки всього є прямих, що проходять через дві вершини куба та є паралельними прямій $AB_1$? \nmtyear{2026}
\vspace{0.2cm}
\begin{enumerate}[label=\textbf{\Alph*}, align=left, leftmargin=2.2em, labelsep=0.6em, itemsep=0.08cm, topsep=0.05cm]
\item жодної
\item лише одна
\item лише дві
\item лише три
\item більше як три
\end{enumerate}
\end{minipage}%
\hfill
\begin{minipage}[c]{0.40\textwidth}
\centering
\begin{tikzpicture}[scale=0.7, thick, line join=round]
    \coordinate (A) at (0,0);
    \coordinate (B) at (1,0.7);
    \coordinate (C) at (3,0.7);
    \coordinate (D) at (2,0);
    \coordinate (A1) at (0,2.4);
    \coordinate (B1) at (1,3.1);
    \coordinate (C1) at (3,3.1);
    \coordinate (D1) at (2,2.4);
    \draw[dashed] (B) -- (A);
    \draw[dashed] (B) -- (C);
    \draw[dashed] (B) -- (B1);
    \draw[thick] (A) -- (D) -- (C) -- (C1) -- (D1) -- (A1) -- cycle;
    \draw[thick] (A1) -- (B1) -- (C1);
    \draw[thick] (D) -- (D1);
    \draw[thick] (A) -- (A1);
    \draw[mainGreen, line width=1.2pt] (A) -- (B1);
    \node[below left] at (A) {$A$};
    \node[below] at (B) {$B$};
    \node[right] at (C) {$C$};
    \node[below right] at (D) {$D$};
    \node[left] at (A1) {$A_1$};
    \node[above left] at (B1) {$B_1$};
    \node[right] at (C1) {$C_1$};
    \node[right] at (D1) {$D_1$};
\end{tikzpicture}
\end{minipage}
\par\vspace{0.3cm}
\solbtn{10}
\noindent
\begin{minipage}[c]{0.56\textwidth}
\zadnum На рисунку зображено куб $ABCDA_1B_1C_1D_1$. Скільки всього є прямих, що проходять через дві вершини куба та є паралельними прямій $AC_1$?
\vspace{0.2cm}
\begin{enumerate}[label=\textbf{\Alph*}, align=left, leftmargin=2.2em, labelsep=0.6em, itemsep=0.08cm, topsep=0.05cm]
\item жодної
\item лише одна
\item лише дві
\item лише три
\item більше як три
\end{enumerate}
\end{minipage}%
\hfill
\begin{minipage}[c]{0.40\textwidth}
\centering
\begin{tikzpicture}[scale=0.9, thick, line join=round]
    \coordinate (A) at (0,0);
    \coordinate (B) at (1,0.7);
    \coordinate (C) at (3,0.7);
    \coordinate (D) at (2,0);
    \coordinate (A1) at (0,2.4);
    \coordinate (B1) at (1,3.1);
    \coordinate (C1) at (3,3.1);
    \coordinate (D1) at (2,2.4);
    \draw[dashed] (B) -- (A);
    \draw[dashed] (B) -- (C);
    \draw[dashed] (B) -- (B1);
    \draw[thick] (A) -- (D) -- (C) -- (C1) -- (D1) -- (A1) -- cycle;
    \draw[thick] (A1) -- (B1) -- (C1);
    \draw[thick] (D) -- (D1);
    \draw[thick] (A) -- (A1);
    \draw[mainGreen, line width=1.2pt] (A) -- (C1);
    \node[below left] at (A) {$A$};
    \node[below] at (B) {$B$};
    \node[right] at (C) {$C$};
    \node[below right] at (D) {$D$};
    \node[left] at (A1) {$A_1$};
    \node[above left] at (B1) {$B_1$};
    \node[right] at (C1) {$C_1$};
    \node[above left] at (D1) {$D_1$};
\end{tikzpicture}
\end{minipage}
\par\vspace{0.3cm}
\noindent
\begin{minipage}[c]{0.56\textwidth}
\zadnum На рисунку зображено куб $ABCDA_1B_1C_1D_1$. Скільки всього є прямих, що проходять через дві вершини куба та є паралельними прямій $AB$?
\vspace{0.2cm}
\begin{enumerate}[label=\textbf{\Alph*}, align=left, leftmargin=2.2em, labelsep=0.6em, itemsep=0.08cm, topsep=0.05cm]
\item жодної
\item лише одна
\item лише дві
\item лише три
\item більше як три
\end{enumerate}
\end{minipage}%
\hfill
\begin{minipage}[c]{0.40\textwidth}
\centering
\begin{tikzpicture}[scale=0.7, thick, line join=round]
    \coordinate (A) at (0,0);
    \coordinate (B) at (1,0.7);
    \coordinate (C) at (3,0.7);
    \coordinate (D) at (2,0);
    \coordinate (A1) at (0,2.4);
    \coordinate (B1) at (1,3.1);
    \coordinate (C1) at (3,3.1);
    \coordinate (D1) at (2,2.4);
    \draw[dashed] (B) -- (C);
    \draw[dashed] (B) -- (B1);
    \draw[mainGreen, line width=1.2pt, dashed] (A) -- (B);
    \draw[thick] (A) -- (D) -- (C) -- (C1) -- (D1) -- (A1) -- cycle;
    \draw[thick] (A1) -- (B1) -- (C1);
    \draw[thick] (D) -- (D1);
    \draw[thick] (A) -- (A1);
    \node[below left] at (A) {$A$};
    \node[below] at (B) {$B$};
    \node[right] at (C) {$C$};
    \node[below right] at (D) {$D$};
    \node[left] at (A1) {$A_1$};
    \node[above left] at (B1) {$B_1$};
    \node[right] at (C1) {$C_1$};
    \node[right] at (D1) {$D_1$};
\end{tikzpicture}
\end{minipage}
\par\vspace{0.3cm}
\taskBlock{Стереометрія: тіла обертання --- сесія 3.06, завдання №6}
\zadnum Конус утворений обертанням \nmtyear{2026}
\vspace{0.1cm}
\begin{enumerate}[label=\textbf{\Alph*}, align=left, leftmargin=2.2em, labelsep=0.6em, itemsep=0.06cm, topsep=0.04cm]
\item рівнобедреного трикутника навколо його основи
\item прямокутного трикутника навколо його гіпотенузи
\item рівностороннього трикутника навколо його сторони
\item прямокутного трикутника навколо його катета
\item прямокутника навколо його сторони
\end{enumerate}
\par\vspace{0.25cm}
\solbtn{13}
\zadnum Циліндр утворений обертанням
\vspace{0.1cm}
\begin{enumerate}[label=\textbf{\Alph*}, align=left, leftmargin=2.2em, labelsep=0.6em, itemsep=0.06cm, topsep=0.04cm]
\item прямокутного трикутника навколо його катета
\item ромба навколо його сторони
\item прямокутника навколо його сторони
\item прямокутника навколо його діагоналі
\item рівнобедреного трикутника навколо його основи
\end{enumerate}
\par\vspace{0.25cm}
\zadnum Унаслідок обертання прямокутного трикутника навколо одного з його катетів утворюється
\vspace{0.1cm}
\begin{enumerate}[label=\textbf{\Alph*}, align=left, leftmargin=2.2em, labelsep=0.6em, itemsep=0.06cm, topsep=0.04cm]
\item циліндр
\item конус
\item куля
\item зрізаний конус
\item два конуси зі спільною основою
\end{enumerate}
\par\vspace{0.25cm}
\taskBlock{Властивості просторових фігур --- сесія 4.06, завдання №5}
\zadtask{Укажіть просторову фігуру, яка не має основи. \nmtyear{2026}
\vspace{0.2cm}
\par\noindent\textbf{А}\ \ призма
\par\noindent\textbf{Б}\ \ конус
\par\noindent\textbf{В}\ \ циліндр
\par\noindent\textbf{Г}\ \ куля
\par\noindent\textbf{Д}\ \ піраміда
\vspace{0.3cm}}
\solbtn{16}
\zadtask{Укажіть просторову фігуру, яка має дві рівні основи, що є кругами.
\vspace{0.2cm}
\par\noindent\textbf{А}\ \ призма
\par\noindent\textbf{Б}\ \ конус
\par\noindent\textbf{В}\ \ циліндр
\par\noindent\textbf{Г}\ \ куля
\par\noindent\textbf{Д}\ \ піраміда
\vspace{0.3cm}}
\zadtask{Укажіть просторову фігуру, бічна поверхня якої складається лише з трикутників.
\vspace{0.2cm}
\par\noindent\textbf{А}\ \ призма
\par\noindent\textbf{Б}\ \ конус
\par\noindent\textbf{В}\ \ циліндр
\par\noindent\textbf{Г}\ \ куля
\par\noindent\textbf{Д}\ \ піраміда
\vspace{0.3cm}}
\taskBlock{Сфера і дотична площина --- сесія 5.06, завдання №5}
\noindent
\begin{minipage}[c]{0.62\textwidth}
\zadnum Площина $\beta$ дотикається сфери в точці $M$ (див. рисунок). Скільки можна провести площин, кожна з яких перпендикулярна до площини $\beta$ і має зі сферою \textit{лише одну} спільну точку? \nmtyear{2026}
\vspace{0.1cm}
\begin{enumerate}[label=\textbf{\Alph*}, align=left, leftmargin=2.2em, labelsep=0.6em, itemsep=0.04cm, topsep=0.03cm]
\item жодної
\item лише одну
\item лише дві
\item лише три
\item безліч
\end{enumerate}
\end{minipage}%
\hfill
\begin{minipage}[c]{0.36\textwidth}
\centering
\begin{tikzpicture}[scale=0.85, thick]
    \draw[fill=yellow!30!orange!20, draw=black, line join=round]
        (-2.5, -0.5) -- (2.5, -0.5) -- (3.8, 1.2) -- (-1.2, 1.2) -- cycle;
    \node at (3.3, 0.9) {$\beta$};
    \coordinate (O) at (0.5, 1.6);
    \def\R{1.6}
    \draw[fill=cyan!15, draw=blue!70!black] (O) circle (\R);
    \draw[draw=blue!70!black, dashed] ($(O) + (-\R, 0)$) arc (180:0:\R cm and 0.4cm);
    \draw[draw=blue!70!black] ($(O) + (-\R, 0)$) arc (180:360:\R cm and 0.4cm);
    \coordinate (M) at (0.5, 0);
    \fill (M) circle (2pt);
    \node[left] at (0.35, 0.05) {$M$};
\end{tikzpicture}
\end{minipage}
\par\vspace{0.2cm}
\solbtn{19}
\noindent
\begin{minipage}[c]{0.62\textwidth}
\zadnum Площина $\beta$ дотикається сфери в точці $M$ (див. рисунок). Скільки можна провести площин, кожна з яких паралельна площині $\beta$ і має зі сферою \textit{лише одну} спільну точку?
\vspace{0.1cm}
\begin{enumerate}[label=\textbf{\Alph*}, align=left, leftmargin=2.2em, labelsep=0.6em, itemsep=0.04cm, topsep=0.03cm]
\item жодної
\item лише одну
\item лише дві
\item лише три
\item безліч
\end{enumerate}
\end{minipage}%
\hfill
\begin{minipage}[c]{0.36\textwidth}
\centering
\begin{tikzpicture}[scale=0.85, thick]
    \draw[fill=yellow!30!orange!20, draw=black, line join=round]
        (-2.5, -0.5) -- (2.5, -0.5) -- (3.8, 1.2) -- (-1.2, 1.2) -- cycle;
    \node at (3.3, 0.9) {$\beta$};
    \coordinate (O) at (0.5, 1.6);
    \def\R{1.6}
    \draw[fill=cyan!15, draw=blue!70!black] (O) circle (\R);
    \draw[draw=blue!70!black, dashed] ($(O) + (-\R, 0)$) arc (180:0:\R cm and 0.4cm);
    \draw[draw=blue!70!black] ($(O) + (-\R, 0)$) arc (180:360:\R cm and 0.4cm);
    \coordinate (M) at (0.5, 0);
    \fill (M) circle (2pt);
    \node[left] at (0.35, 0.05) {$M$};
\end{tikzpicture}
\end{minipage}
\par\vspace{0.2cm}
\noindent
\begin{minipage}[c]{0.62\textwidth}
\zadnum Площина $\beta$ дотикається сфери в точці $M$ (див. рисунок). Скільки можна провести прямих, кожна з яких лежить у площині $\beta$ і має зі сферою \textit{лише одну} спільну точку?
\vspace{0.1cm}
\begin{enumerate}[label=\textbf{\Alph*}, align=left, leftmargin=2.2em, labelsep=0.6em, itemsep=0.04cm, topsep=0.03cm]
\item жодної
\item лише одну
\item лише дві
\item лише три
\item безліч
\end{enumerate}
\end{minipage}%
\hfill
\begin{minipage}[c]{0.36\textwidth}
\centering
\begin{tikzpicture}[scale=0.85, thick]
    \draw[fill=yellow!30!orange!20, draw=black, line join=round]
        (-2.5, -0.5) -- (2.5, -0.5) -- (3.8, 1.2) -- (-1.2, 1.2) -- cycle;
    \node at (3.3, 0.9) {$\beta$};
    \coordinate (O) at (0.5, 1.6);
    \def\R{1.6}
    \draw[fill=cyan!15, draw=blue!70!black] (O) circle (\R);
    \draw[draw=blue!70!black, dashed] ($(O) + (-\R, 0)$) arc (180:0:\R cm and 0.4cm);
    \draw[draw=blue!70!black] ($(O) + (-\R, 0)$) arc (180:360:\R cm and 0.4cm);
    \coordinate (M) at (0.5, 0);
    \fill (M) circle (2pt);
    \node[left] at (0.35, 0.05) {$M$};
\end{tikzpicture}
\end{minipage}
\par\vspace{0.2cm}
\subTypeTitle{Вектори і координати в просторі}
\taskBlock{Модуль вектора у просторі --- сесія 23.05, завдання №8}
\zadtask{Визначте довжину (модуль) вектора $\vec{a}(5;\, 2;\, -1)$. \nmtyear{2026}}
\answerTable{$\sqrt{28}$}{$6$}{$\sqrt{30}$}{$\sqrt{32}$}{$15$}
\solbtn{22}
\zadtask{Визначте довжину (модуль) вектора $\vec{b}(-3;\, 6;\, 2)$.}
\answerTable{$5$}{$\sqrt{47}$}{$11$}{$7$}{$\sqrt{41}$}
\zadtask{Дано точки $A(1;\, -2;\, 3)$ та $B(3;\, 2;\, 1)$. Визначте довжину вектора $\vec{AB}$.}
\answerTable{$2\sqrt{5}$}{$6$}{$2\sqrt{6}$}{$\sqrt{26}$}{$\sqrt{14}$}
\taskBlock{Стереометрія: сфера в координатах --- сесія 30.05, завдання №10}
\zadtask{У прямокутній системі координат у просторі задано сферу та дві точки $M(-4;\, 5;\, 0)$ і $N(4;\, -3;\, 4)$, які належать цій сфері, причому відстань між ними є \textit{найбільшою} серед усіх можливих відстаней між точками цієї сфери. Знайдіть радіус сфери. \nmtyear{2026}}
\answerTable{$4$}{$6$}{$8$}{$12$}{$24$}
\solbtn{25}
\zadtask{У прямокутній системі координат у просторі задано сферу та дві точки $M(3;\, -2;\, 4)$ і $N(3;\, 4;\, -4)$, які належать цій сфері, причому відстань між ними є \textit{найбільшою} серед усіх можливих відстаней між точками цієї сфери. Знайдіть радіус сфери.}
\answerTable{$2{,}5$}{$5$}{$7$}{$10$}{$20$}
\zadtask{У прямокутній системі координат у просторі сферу задано рівнянням $(x-3)^2 + (y+4)^2 + (z-12)^2 = 100$. Знайдіть радіус цієї сфери.}
\answerTable{$5$}{$8$}{$10$}{$20$}{$40$}
\taskBlock{Відстань у просторі від початку координат --- сесія 1.06, завдання №15}
\zadtask{У прямокутній системі координат у просторі задано точку $O(0;\,0;\,0)$. Укажіть з-поміж наведених точку, відстань від якої до точки $O$ \textit{не дорівнює} $10$. \nmtyear{2026}}
\answerTable{$(0;\,6;\,8)$}{$(10;\,0;\,0)$}{$(9;\,1;\,0)$}{$(-6;\,0;\,8)$}{$(0;\,-8;\,-6)$}
\solbtn{28}
\zadtask{У прямокутній системі координат у просторі задано точку $O(0;\,0;\,0)$. Укажіть з-поміж наведених точку, відстань від якої до точки $O$ \textit{не дорівнює} $13$.}
\answerTable{$(12;\,4;\,3)$}{$(13;\,0;\,0)$}{$(0;\,5;\,12)$}{$(8;\,2;\,5)$}{$(3;\,-12;\,-4)$}
\zadtask{У прямокутній системі координат у просторі задано точки $O(0;\,0;\,0)$, $A(1;\,2;\,2)$, $B(2;\,3;\,6)$, $C(6;\,6;\,7)$, $D(4;\,4;\,2)$, $E(0;\,3;\,4)$. Укажіть точку, \textit{найбільш віддалену} від $O$.}
\answerTable{$A$}{$D$}{$C$}{$B$}{$E$}
\taskBlock{Вектор у просторі: координати точки --- сесія 2.06, завдання №12}
\zadtask{У прямокутній системі координат у просторі задано вектор $\vec{a}\,(5;\, -3;\, 8)$ та точки $K$ і $M(-7;\, 1;\, 4)$. Знайдіть координати точки $K$, якщо $\vec{a} = \overrightarrow{KM}$. \nmtyear{2026}}
\answerTable{$(-2;\, -2;\, 12)$}{$(-12;\, 4;\, -4)$}{$(-2;\, 4;\, 4)$}{$(2;\, 2;\, -12)$}{$(12;\, -4;\, 4)$}
\solbtn{31}
\zadtask{У прямокутній системі координат у просторі задано вектор $\vec{a}\,(3;\, -7;\, 2)$ та точки $K$ і $M(-4;\, 5;\, 1)$. Знайдіть координати точки $K$, якщо $\vec{a} = \overrightarrow{KM}$.}
\answerTable{$(-1;\, -2;\, 3)$}{$(7;\, -12;\, 1)$}{$(-7;\, 12;\, -1)$}{$(1;\, 2;\, -3)$}{$(-7;\, -2;\, -1)$}
\zadtask{У прямокутній системі координат у просторі задано вектор $\vec{a}\,(-2;\, 6;\, -5)$ та точки $K(3;\, -1;\, 4)$ і $M$. Знайдіть координати точки $M$, якщо $\vec{a} = \overrightarrow{MK}$.}
\answerTable{$(5;\, -7;\, 9)$}{$(1;\, 5;\, -1)$}{$(-5;\, 7;\, -9)$}{$(-1;\, -5;\, 1)$}{$(1;\, -5;\, -1)$}
\taskBlock{Стереометрія: вектори у просторі --- сесія 3.06, завдання №12}
\zadtask{У прямокутній системі координат у просторі задано вектори $\vec{a}\,(-2;\, 2;\, 1)$ і $\vec{b}\,(4;\, -3;\, 0)$. Визначте модуль вектора $\vec{c} = \vec{a} + \vec{b}$. \nmtyear{2026}}
\answerTable{$\sqrt{5}$}{$2$}{$\sqrt{62}$}{$\sqrt{6}$}{$14$}
\solbtn{34}
\zadtask{У прямокутній системі координат у просторі задано вектори $\vec{a}\,(3;\, -1;\, 2)$ і $\vec{b}\,(-1;\, 4;\, 2)$. Визначте модуль вектора $\vec{c} = \vec{a} + \vec{b}$.}
\answerTable{$3$}{$\sqrt{29}$}{$\sqrt{21}$}{$29$}{$\sqrt{14}$}
\zadtask{У прямокутній системі координат у просторі задано вектори $\vec{a}\,(2;\, -3;\, 1)$ і $\vec{b}\,(4;\, -1;\, -1)$. Визначте модуль вектора $\vec{c} = \vec{a} - \vec{b}$.}
\answerTable{$2\sqrt{3}$}{$2\sqrt{13}$}{$12$}{$\sqrt{6}$}{$4$}
\taskBlock{Циліндр у прямокутній системі координат --- сесія 4.06, завдання №12}
\zadtask{У прямокутній системі координат у просторі задано циліндр та точки $O_1(-3;\, 2;\, 1)$ і $O_2(3;\, 4;\, 4)$, що є центрами його основ. Визначте довжину твірної циліндра. \nmtyear{2026}}
\answerTable{$7$}{$3\sqrt{5}$}{$\sqrt{13}$}{$\sqrt{61}$}{$3$}
\solbtn{37}
\zadtask{У прямокутній системі координат у просторі задано циліндр та точки $O_1(1;\, -2;\, 2)$ і $O_2(5;\, 2;\, 0)$, що є центрами його основ. Визначте довжину твірної циліндра.}
\answerTable{$4$}{$6$}{$2\sqrt{5}$}{$\sqrt{34}$}{$2\sqrt{2}$}
\zadtask{У прямокутній системі координат у просторі задано циліндр та точки $O_1(2;\, 0;\, -1)$ і $O_2(6;\, 4;\, 11)$, що є центрами його основ. Визначте координати центра циліндра (середини його осі).}
\answerTable{$(4;\, 2;\, 5)$}{$(8;\, 4;\, 10)$}{$(2;\, 2;\, 6)$}{$(4;\, 2;\, 6)$}{$(3;\, 2;\, 5)$}
\taskBlock{Довжина вектора у просторі --- сесія 5.06, завдання №7}
\zadtask{У прямокутній системі координат у просторі задано точки $A(1;\,2;\,3)$ і $B(3;\,1;\,1)$. Визначте довжину вектора $\overrightarrow{AB}$. \nmtyear{2026}}
\answerTable{$3$}{$\sqrt{6}$}{$9$}{$\sqrt{14}$}{$2\sqrt{3}$}
\solbtn{40}
\zadtask{У прямокутній системі координат у просторі задано точки $A(1;\,0;\,2)$ і $B(3;\,3;\,8)$. Визначте довжину вектора $\overrightarrow{AB}$.}
\answerTable{$\sqrt{11}$}{$7$}{$49$}{$\sqrt{29}$}{$5$}
\zadtask{У прямокутній системі координат у просторі задано точки $A(1;\,2;\,3)$, $B(4;\,4;\,5)$ і $C(3;\,5;\,9)$. Визначте довжину вектора $\overrightarrow{AB} + \overrightarrow{BC}$.}
\answerTable{$\sqrt{17}$}{$7$}{$49$}{$5$}{$\sqrt{61}$}
\taskBlock{Середина відрізка у просторі --- сесія 5.06, завдання №12}
\zadtask{Визначте координати середини відрізка $AB$, якщо $A(1;\,2;\,3)$ і $B(3;\,1;\,1)$.}
\answerTable{$(-2;\,-1;\,-2)$}{$(1;\,-0{,}5;\,-1)$}{$(2;\,1{,}5;\,2)$}{$(4;\,3;\,1)$}{$(-1;\,-0{,}5;\,-1)$}
\solbtn{43}
\zadtask{Визначте координати середини відрізка $AB$, якщо $A(2;\,-3;\,5)$ і $B(4;\,1;\,-1)$.}
\answerTable{$(6;\,-2;\,4)$}{$(1;\,2;\,-3)$}{$(-1;\,-2;\,3)$}{$(3;\,-1;\,2)$}{$(3;\,-2;\,2)$}
\zadtask{Точка $A'$ симетрична точці $A(1;\,4;\,2)$ відносно точки $O(3;\,-1;\,5)$. Визначте координати точки $A'$.}
\answerTable{$(5;\,-6;\,8)$}{$(2;\,1{,}5;\,3{,}5)$}{$(-1;\,9;\,-1)$}{$(4;\,3;\,7)$}{$(2;\,-5;\,3)$}
\typeTitle{Завдання з короткою відповіддю}
\subTypeTitle{Стереометричні розрахунки}
\taskBlock{Циліндр і конус зі спільною віссю --- сесія 23.05, завдання №19}
\noindent
\begin{minipage}[c]{0.60\textwidth}
\zadnum На рисунку зображено циліндр і конус, що мають спільну вісь. Радіус основи конуса вдвічі менший за радіус циліндра. Висота конуса дорівнює $15$~см, а твірна конуса~--- $17$~см. Знайдіть об'єм циліндра (у см$^3$), якщо площа його бічної поверхні дорівнює $640\pi$~см$^2$. У відповіді запишіть значення $\dfrac{V}{\pi}$. \nmtyear{2026}
\end{minipage}%
\hfill
\begin{minipage}[c]{0.36\textwidth}
\centering
\begin{tikzpicture}[scale=0.7, thick, line join=round]
    \draw[dashed] (-1.1,0) arc (180:0:1.1 and 0.35);
    \draw (-1.1,0) arc (180:360:1.1 and 0.35);
    \draw (-1.1,0) -- (-1.1,3.2);
    \draw (1.1,0) -- (1.1,3.2);
    \draw (0,3.2) ellipse (1.1 and 0.35);
    \begin{scope}[shift={(3.7,0)}]
        \draw[dashed] (-0.8,0) arc (180:0:0.8 and 0.25);
        \draw (-0.8,0) arc (180:360:0.8 and 0.25);
        \draw (-0.8,0) -- (0,3.5);
        \draw (0.8,0) -- (0,3.5);
    \end{scope}
\end{tikzpicture}
\end{minipage}
\par\vspace{0.2cm}
\nmtAnswerBox
\solbtn{46}
\vspace{0.3cm}
\noindent
\begin{minipage}[c]{0.60\textwidth}
\zadnum На рисунку зображено циліндр і конус, що мають спільну вісь. Радіус основи конуса вдвічі менший за радіус циліндра, а висоти циліндра і конуса рівні. Висота конуса дорівнює $12$~см, а його твірна утворює з площиною основи кут $45^\circ$. Знайдіть об'єм циліндра (у см$^3$). У відповіді запишіть значення $\dfrac{V}{\pi}$.
\end{minipage}%
\hfill
\begin{minipage}[c]{0.36\textwidth}
\centering
\begin{tikzpicture}[scale=0.7, thick, line join=round]
    \draw[dashed] (-1.1,0) arc (180:0:1.1 and 0.35);
    \draw (-1.1,0) arc (180:360:1.1 and 0.35);
    \draw (-1.1,0) -- (-1.1,3.0);
    \draw (1.1,0) -- (1.1,3.0);
    \draw (0,3.0) ellipse (1.1 and 0.35);
    \begin{scope}[shift={(3.7,0)}]
        \draw[dashed] (-0.8,0) arc (180:0:0.8 and 0.25);
        \draw (-0.8,0) arc (180:360:0.8 and 0.25);
        \draw (-0.8,0) -- (0,3.0);
        \draw (0.8,0) -- (0,3.0);
    \end{scope}
\end{tikzpicture}
\end{minipage}
\par\vspace{0.2cm}
\nmtAnswerBox
\vspace{0.3cm}
\noindent
\begin{minipage}[c]{0.60\textwidth}
\zadnum На рисунку зображено циліндр і конус, що мають спільну вісь та \textit{рівні радіуси основ}, які дорівнюють $2\sqrt{3}$~см. Твірна конуса утворює з площиною основи кут $60^\circ$, а висота циліндра вдвічі менша за висоту конуса. Знайдіть об'єм циліндра (у см$^3$). У відповіді запишіть значення $\dfrac{V}{\pi}$.
\end{minipage}%
\hfill
\begin{minipage}[c]{0.36\textwidth}
\centering
\begin{tikzpicture}[scale=0.7, thick, line join=round]
    \draw[dashed] (-1.1,0) arc (180:0:1.1 and 0.35);
    \draw (-1.1,0) arc (180:360:1.1 and 0.35);
    \draw (-1.1,0) -- (-1.1,1.8);
    \draw (1.1,0) -- (1.1,1.8);
    \draw (0,1.8) ellipse (1.1 and 0.35);
    \begin{scope}[shift={(3.7,0)}]
        \draw[dashed] (-1.1,0) arc (180:0:1.1 and 0.35);
        \draw (-1.1,0) arc (180:360:1.1 and 0.35);
        \draw (-1.1,0) -- (0,3.5);
        \draw (1.1,0) -- (0,3.5);
    \end{scope}
\end{tikzpicture}
\end{minipage}
\par\vspace{0.2cm}
\nmtAnswerBox
\taskBlock{Стереометрія: призма і конус --- сесія 30.05, завдання №21}
\noindent
\begin{minipage}[c]{0.58\textwidth}
\zadnum Пряма чотирикутна призма й конус мають рівні висоти. В основі призми лежить рівнобічна трапеція, у яку можна вписати коло. Висота трапеції дорівнює діаметру основи конуса. Бічна сторона трапеції дорівнює $26$ см, а площа основи призми~--- $624$ см$^2$. Твірна конуса дорівнює $20$ см. Визначте площу (у см$^2$) бічної поверхні призми. \nmtyear{2026}
\end{minipage}%
\hfill
\begin{minipage}[c]{0.40\textwidth}
\centering
\begin{tikzpicture}[scale=0.5, thick, line join=round]
    \coordinate (A) at (0,0);
    \coordinate (B) at (5,0);
    \coordinate (C) at (4,1);
    \coordinate (D) at (1,1);
    \coordinate (A1) at (0,3);
    \coordinate (B1) at (5,3);
    \coordinate (C1) at (4,4);
    \coordinate (D1) at (1,4);
    \draw[thick] (A) -- (A1);
    \draw[thick] (B) -- (B1);
    \draw[thick] (C) -- (C1);
    \draw[thick, dashed] (D) -- (D1);
    \draw[thick] (A1) -- (B1) -- (C1) -- (D1) -- cycle;
    \draw[thick] (A) -- (B);
    \draw[thick] (B) -- (C);
    \draw[thick, dashed] (C) -- (D);
    \draw[thick, dashed] (D) -- (A);
    \coordinate (V) at (9, 4);
    \draw[thick, dashed] (7.2, 0) arc (180:0:1.8 and 0.5);
    \draw[thick] (7.2, 0) arc (180:360:1.8 and 0.5);
    \draw[thick] (V) -- (7.2, 0);
    \draw[thick] (V) -- (10.8, 0);
    \draw[thick, dashed] (V) -- (9, 0);
\end{tikzpicture}
\end{minipage}
\par\vspace{0.2cm}
\nmtAnswerBox
\solbtn{49}
\zadtask{Пряма чотирикутна призма й конус мають рівні висоти. В основі призми лежить рівнобічна трапеція, у яку можна вписати коло; гострий кут трапеції дорівнює $30^\circ$, а бічна сторона~--- $24$ см. Висота трапеції дорівнює діаметру основи конуса, а твірна конуса утворює з площиною його основи кут $45^\circ$. Визначте площу (у см$^2$) бічної поверхні призми.}
\nmtAnswerBox
\zadtask{Пряма чотирикутна призма й конус мають рівні висоти. В основі призми лежить рівнобічна трапеція, у яку можна вписати коло. Висота трапеції дорівнює діаметру основи конуса. Бічна сторона трапеції дорівнює $10$ см, а площа основи призми~--- $80$ см$^2$. Твірна конуса дорівнює $5$ см. Визначте \textit{об'єм} призми (у см$^3$).}
\nmtAnswerBox
\taskBlock{Конус і правильна трикутна піраміда, об'єм конуса --- сесія 1.06, завдання №21}
\noindent
\begin{minipage}[c]{0.62\textwidth}
\zadnum Конус і правильна трикутна піраміда мають \textit{рівні висоти}. Радіус основи конуса дорівнює радіусу кола, \textit{вписаного} в основу піраміди, і дорівнює $4\sqrt{3}$. Площа бічної поверхні піраміди дорівнює $288$. Знайдіть об'єм конуса. У відповідь запишіть значення $\dfrac{V}{\pi}$. \nmtyear{2026}
\end{minipage}%
\hfill
\begin{minipage}[c]{0.36\textwidth}
\centering
\begin{tikzpicture}[scale=0.85, thick, line join=round]
    \draw[thick, dashed] (-1.3,0) arc (180:0:1.3 and 0.38);
    \draw[thick] (-1.3,0) arc (180:360:1.3 and 0.38);
    \coordinate (CV) at (0,3);
    \draw[thick] (CV) -- (-1.3,0);
    \draw[thick] (CV) -- (1.3,0);
    \draw[thick, dashed] (CV) -- (0,0);
    \draw[<-, >=stealth, gray] (0,0) -- (1.3,0);
    \node[below, font=\scriptsize] at (0.7,-0.05) {$R$};
    \coordinate (A) at (3.2,-0.35);
    \coordinate (B) at (5.6,-0.35);
    \coordinate (C) at (4.4,0.45);
    \coordinate (S) at (4.4,3);
    \draw[thick] (A) -- (B) -- (S) -- cycle;
    \draw[thick] (A) -- (S);
    \draw[thick, dashed] (A) -- (C) -- (B);
    \draw[thick, dashed] (C) -- (S);
\end{tikzpicture}
\end{minipage}
\nmtAnswerBox
\solbtn{52}
\noindent
\begin{minipage}[c]{0.62\textwidth}
\zadnum Конус і правильна трикутна піраміда мають \textit{рівні висоти}. Радіус основи конуса дорівнює радіусу кола, \textit{вписаного} в основу піраміди. Сторона основи піраміди дорівнює $12$, а апофема піраміди утворює з площиною основи кут $60^\circ$. Знайдіть об'єм конуса. У відповідь запишіть значення $\dfrac{V}{\pi}$.
\end{minipage}%
\hfill
\begin{minipage}[c]{0.36\textwidth}
\centering
\begin{tikzpicture}[scale=0.85, thick, line join=round]
    \draw[thick, dashed] (-1.3,0) arc (180:0:1.3 and 0.38);
    \draw[thick] (-1.3,0) arc (180:360:1.3 and 0.38);
    \coordinate (CV) at (0,3);
    \draw[thick] (CV) -- (-1.3,0);
    \draw[thick] (CV) -- (1.3,0);
    \draw[thick, dashed] (CV) -- (0,0);
    \draw[<-, >=stealth, gray] (0,0) -- (1.3,0);
    \node[below, font=\scriptsize] at (0.7,-0.05) {$R$};
    \coordinate (A) at (3.2,-0.35);
    \coordinate (B) at (5.6,-0.35);
    \coordinate (C) at (4.4,0.45);
    \coordinate (S) at (4.4,3);
    \draw[thick] (A) -- (B) -- (S) -- cycle;
    \draw[thick] (A) -- (S);
    \draw[thick, dashed] (A) -- (C) -- (B);
    \draw[thick, dashed] (C) -- (S);
\end{tikzpicture}
\end{minipage}
\nmtAnswerBox
\noindent
\begin{minipage}[c]{0.62\textwidth}
\zadnum Конус і правильна трикутна піраміда мають \textit{рівні висоти}. Радіус основи конуса дорівнює радіусу кола, \textit{описаного} навколо основи піраміди. Радіус кола, \textit{вписаного} в основу піраміди, дорівнює $3\sqrt{3}$, а площа бічної поверхні піраміди дорівнює $162$. Знайдіть об'єм конуса. У відповідь запишіть значення $\dfrac{V}{\pi}$.
\end{minipage}%
\hfill
\begin{minipage}[c]{0.36\textwidth}
\centering
\begin{tikzpicture}[scale=0.85, thick, line join=round]
    \draw[thick, dashed] (-1.3,0) arc (180:0:1.3 and 0.38);
    \draw[thick] (-1.3,0) arc (180:360:1.3 and 0.38);
    \coordinate (CV) at (0,3);
    \draw[thick] (CV) -- (-1.3,0);
    \draw[thick] (CV) -- (1.3,0);
    \draw[thick, dashed] (CV) -- (0,0);
    \draw[<-, >=stealth, gray] (0,0) -- (1.3,0);
    \node[below, font=\scriptsize] at (0.7,-0.05) {$R$};
    \coordinate (A) at (3.2,-0.35);
    \coordinate (B) at (5.6,-0.35);
    \coordinate (C) at (4.4,0.45);
    \coordinate (S) at (4.4,3);
    \draw[thick] (A) -- (B) -- (S) -- cycle;
    \draw[thick] (A) -- (S);
    \draw[thick, dashed] (A) -- (C) -- (B);
    \draw[thick, dashed] (C) -- (S);
\end{tikzpicture}
\end{minipage}
\nmtAnswerBox
\taskBlock{Об'єми півсфери та конуса --- сесія 2.06, завдання №21}
\noindent
\begin{minipage}[c]{0.62\textwidth}
\zadnum Радіус основи й твірна конуса відносяться як $3 : 5$, а площа його бічної поверхні дорівнює $135\pi$. Площа повної поверхні півсфери дорівнює $243\pi$. У скільки разів об'єм півсфери більший за об'єм конуса? \nmtyear{2026}
\end{minipage}%
\hfill
\begin{minipage}[c]{0.36\textwidth}
\centering
\begin{tikzpicture}[scale=0.8, thick, line join=round]
    \draw[thick, dashed] (-1.0,0) arc (180:0:1.0 and 0.28);
    \draw[thick] (-1.0,0) arc (180:360:1.0 and 0.28);
    \draw[thick] (-1.0,0) -- (0,2.4) -- (1.0,0);
    \begin{scope}[shift={(3.4,0)}]
        \draw[thick] (-1.1,0) arc (180:0:1.1 and 1.1);
        \draw[thick] (-1.1,0) -- (1.1,0);
        \draw[thick, dashed] (-1.1,0) arc (180:0:1.1 and 0.28);
        \draw[thick] (-1.1,0) arc (180:360:1.1 and 0.28);
    \end{scope}
\end{tikzpicture}
\end{minipage}
\nmtAnswerBox
\solbtn{55}
\noindent
\begin{minipage}[c]{0.62\textwidth}
\zadnum Конус і півсфера мають \textit{рівні радіуси основ}. Твірна конуса утворює з площиною його основи кут $45^\circ$. У скільки разів об'єм півсфери більший за об'єм конуса?
\end{minipage}%
\hfill
\begin{minipage}[c]{0.36\textwidth}
\centering
\begin{tikzpicture}[scale=0.8, thick, line join=round]
    \draw[thick, dashed] (-1.1,0) arc (180:0:1.1 and 0.28);
    \draw[thick] (-1.1,0) arc (180:360:1.1 and 0.28);
    \draw[thick] (-1.1,0) -- (0,1.1) -- (1.1,0);
    \begin{scope}[shift={(3.4,0)}]
        \draw[thick] (-1.1,0) arc (180:360:1.1 and 1.1);
        \draw[thick] (-1.1,0) -- (1.1,0);
        \draw[thick, dashed] (-1.1,0) arc (180:0:1.1 and 0.28);
        \draw[thick] (-1.1,0) arc (180:360:1.1 and 0.28);
    \end{scope}
\end{tikzpicture}
\end{minipage}
\nmtAnswerBox
\noindent
\begin{minipage}[c]{0.62\textwidth}
\zadnum Радіус основи конуса дорівнює $6$ см, а висота~--- $16$ см. Діаметр півсфери дорівнює $24$ см. У скільки разів об'єм півсфери більший за об'єм конуса?
\end{minipage}%
\hfill
\begin{minipage}[c]{0.36\textwidth}
\centering
\begin{tikzpicture}[scale=0.8, thick, line join=round]
    \draw[thick, dashed] (-0.8,0) arc (180:0:0.8 and 0.22);
    \draw[thick] (-0.8,0) arc (180:360:0.8 and 0.22);
    \draw[thick] (-0.8,0) -- (0,2.4) -- (0.8,0);
    \begin{scope}[shift={(3.4,0)}]
        \draw[thick] (-1.3,0) arc (180:360:1.3 and 1.3);
        \draw[thick] (-1.3,0) -- (1.3,0);
        \draw[thick, dashed] (-1.3,0) arc (180:0:1.3 and 0.32);
        \draw[thick] (-1.3,0) arc (180:360:1.3 and 0.32);
    \end{scope}
\end{tikzpicture}
\end{minipage}
\nmtAnswerBox
\taskBlock{Стереометрія: циліндр і паралелепіпед --- сесія 3.06, завдання №21}
\noindent
\begin{minipage}[c]{0.62\textwidth}
\zadnum На рисунку зображено циліндр і прямокутний паралелепіпед. Діаметр основи циліндра дорівнює діагоналі основи паралелепіпеда. Довжина діагоналі паралелепіпеда дорівнює $39$ см, сторони основи~--- $9$ см і $12$ см. Висота циліндра відноситься до висоти паралелепіпеда як $2 : 3$. Визначте об'єм $V$ циліндра (у см$^3$). У відповідь запишіть значення $\dfrac{V}{\pi}$. \nmtyear{2026}
\end{minipage}%
\hfill
\begin{minipage}[c]{0.36\textwidth}
\centering
\begin{tikzpicture}[scale=0.55, thick, line join=round]
    \draw[thick, dashed] (-0.9,0) arc (180:0:0.9 and 0.26);
    \draw[thick] (-0.9,0) arc (180:360:0.9 and 0.26);
    \draw[thick] (-0.9,0) -- (-0.9,4.0);
    \draw[thick] (0.9,0) -- (0.9,4.0);
    \draw[thick] (0,4.0) ellipse (0.9 and 0.26);
    \begin{scope}[shift={(3.2,0)}]
        \coordinate (A) at (0,0);
        \coordinate (B) at (2,0);
        \coordinate (C) at (2.7,0.7);
        \coordinate (D) at (0.7,0.7);
        \draw[thick] (A) -- (B) -- (C);
        \draw[thick, dashed] (C) -- (D) -- (A);
        \draw[thick] (A) -- (0,2.6);
        \draw[thick] (B) -- (2,2.6);
        \draw[thick] (C) -- (2.7,3.3);
        \draw[thick, dashed] (D) -- (0.7,3.3);
        \draw[thick] (0,2.6) -- (2,2.6) -- (2.7,3.3) -- (0.7,3.3) -- cycle;
    \end{scope}
\end{tikzpicture}
\end{minipage}
\nmtAnswerBox
\solbtn{58}
\noindent
\begin{minipage}[c]{0.62\textwidth}
\zadnum На рисунку зображено циліндр і прямокутний паралелепіпед. Діаметр основи циліндра дорівнює діагоналі основи паралелепіпеда, а висота циліндра дорівнює висоті паралелепіпеда. Сторони основи паралелепіпеда дорівнюють $6$ см і $8$ см, а діагональ паралелепіпеда утворює з площиною основи кут $45^\circ$. Визначте об'єм $V$ циліндра (у см$^3$). У відповідь запишіть значення $\dfrac{V}{\pi}$.
\end{minipage}%
\hfill
\begin{minipage}[c]{0.36\textwidth}
\centering
\begin{tikzpicture}[scale=0.55, thick, line join=round]
    \draw[thick, dashed] (-0.9,0) arc (180:0:0.9 and 0.26);
    \draw[thick] (-0.9,0) arc (180:360:0.9 and 0.26);
    \draw[thick] (-0.9,0) -- (-0.9,2.6);
    \draw[thick] (0.9,0) -- (0.9,2.6);
    \draw[thick] (0,2.6) ellipse (0.9 and 0.26);
    \begin{scope}[shift={(3.2,0)}]
        \coordinate (A) at (0,0);
        \coordinate (B) at (2,0);
        \coordinate (C) at (2.7,0.7);
        \coordinate (D) at (0.7,0.7);
        \draw[thick] (A) -- (B) -- (C);
        \draw[thick, dashed] (C) -- (D) -- (A);
        \draw[thick] (A) -- (0,2.6);
        \draw[thick] (B) -- (2,2.6);
        \draw[thick] (C) -- (2.7,3.3);
        \draw[thick, dashed] (D) -- (0.7,3.3);
        \draw[thick] (0,2.6) -- (2,2.6) -- (2.7,3.3) -- (0.7,3.3) -- cycle;
    \end{scope}
\end{tikzpicture}
\end{minipage}
\nmtAnswerBox
\noindent
\begin{minipage}[c]{0.62\textwidth}
\zadnum На рисунку зображено циліндр і прямокутний паралелепіпед. Діаметр основи циліндра дорівнює діагоналі основи паралелепіпеда, а \textit{висота циліндра дорівнює діагоналі паралелепіпеда}. Сторони основи паралелепіпеда~--- $12$ см і $16$ см, висота паралелепіпеда~--- $15$ см. Визначте об'єм $V$ циліндра (у см$^3$). У відповідь запишіть значення $\dfrac{V}{\pi}$.
\end{minipage}%
\hfill
\begin{minipage}[c]{0.36\textwidth}
\centering
\begin{tikzpicture}[scale=0.55, thick, line join=round]
    \draw[thick, dashed] (-0.9,0) arc (180:0:0.9 and 0.26);
    \draw[thick] (-0.9,0) arc (180:360:0.9 and 0.26);
    \draw[thick] (-0.9,0) -- (-0.9,2.6);
    \draw[thick] (0.9,0) -- (0.9,2.6);
    \draw[thick] (0,2.6) ellipse (0.9 and 0.26);
    \begin{scope}[shift={(3.2,0)}]
        \coordinate (A) at (0,0);
        \coordinate (B) at (2,0);
        \coordinate (C) at (2.7,0.7);
        \coordinate (D) at (0.7,0.7);
        \draw[thick] (A) -- (B) -- (C);
        \draw[thick, dashed] (C) -- (D) -- (A);
        \draw[thick] (A) -- (0,2.6);
        \draw[thick] (B) -- (2,2.6);
        \draw[thick] (C) -- (2.7,3.3);
        \draw[thick, dashed] (D) -- (0.7,3.3);
        \draw[thick] (0,2.6) -- (2,2.6) -- (2.7,3.3) -- (0.7,3.3) -- cycle;
    \end{scope}
\end{tikzpicture}
\end{minipage}
\nmtAnswerBox
\taskBlock{Об'єм призми (призма і піраміда) --- сесія 4.06, завдання №21}
\noindent
\begin{minipage}[c]{0.58\textwidth}
\zadnum На рисунку зображено правильну чотирикутну призму та правильну чотирикутну піраміду. Сторона основи піраміди вдвічі менша від діагоналі основи призми. Усі ребра піраміди рівні й дорівнюють $3\sqrt{2}$~см. Знайдіть об'єм (у см$^3$) призми, якщо сума довжин висот призми та піраміди дорівнює $13$~см. \nmtyear{2026}
\par\vspace{0.4cm}
\nmtAnswerBox
\end{minipage}%
\hfill
\begin{minipage}[c]{0.40\textwidth}
\centering
\begin{tikzpicture}[scale=0.65, thick, line join=round]
    \coordinate (A) at (0,0);
    \coordinate (B) at (1.8,0);
    \coordinate (C) at (2.6, 0.7);
    \coordinate (D) at (0.8, 0.7);
    \coordinate (A1) at (0,3.8);
    \coordinate (B1) at (1.8,3.8);
    \coordinate (C1) at (2.6, 4.5);
    \coordinate (D1) at (0.8, 4.5);
    \draw (A) -- (B) -- (C) -- (C1) -- (D1) -- (A1) -- cycle;
    \draw (B) -- (B1) -- (C1);
    \draw (A1) -- (B1);
    \draw[dashed] (A) -- (D) -- (C);
    \draw[dashed] (D) -- (D1);
    \begin{scope}[shift={(3.8,0)}]
        \coordinate (P_A) at (0,0);
        \coordinate (P_B) at (2.2,0);
        \coordinate (P_C) at (3.2, 0.9);
        \coordinate (P_D) at (1.0, 0.9);
        \coordinate (P_S) at (1.6, 4.0);
        \draw (P_A) -- (P_B) -- (P_C) -- (P_S) -- cycle;
        \draw (P_B) -- (P_S);
        \draw[dashed] (P_A) -- (P_D) -- (P_C);
        \draw[dashed] (P_D) -- (P_S);
    \end{scope}
\end{tikzpicture}
\end{minipage}
\par\vspace{0.3cm}
\solbtn{61}
\noindent
\begin{minipage}[c]{0.58\textwidth}
\zadnum На рисунку зображено правильну чотирикутну призму та правильну чотирикутну піраміду. Сторона основи піраміди вдвічі менша від діагоналі основи призми і дорівнює $6$~см. Апофема піраміди утворює з площиною її основи кут $45^\circ$. Знайдіть об'єм (у см$^3$) призми, якщо сума довжин висот призми та піраміди дорівнює $13$~см.
\par\vspace{0.4cm}
\nmtAnswerBox
\end{minipage}%
\hfill
\begin{minipage}[c]{0.40\textwidth}
\centering
\begin{tikzpicture}[scale=0.65, thick, line join=round]
    \coordinate (A) at (0,0);
    \coordinate (B) at (1.8,0);
    \coordinate (C) at (2.6, 0.7);
    \coordinate (D) at (0.8, 0.7);
    \coordinate (A1) at (0,3.8);
    \coordinate (B1) at (1.8,3.8);
    \coordinate (C1) at (2.6, 4.5);
    \coordinate (D1) at (0.8, 4.5);
    \draw (A) -- (B) -- (C) -- (C1) -- (D1) -- (A1) -- cycle;
    \draw (B) -- (B1) -- (C1);
    \draw (A1) -- (B1);
    \draw[dashed] (A) -- (D) -- (C);
    \draw[dashed] (D) -- (D1);
    \begin{scope}[shift={(3.8,0)}]
        \coordinate (P_A) at (0,0);
        \coordinate (P_B) at (2.2,0);
        \coordinate (P_C) at (3.2, 0.9);
        \coordinate (P_D) at (1.0, 0.9);
        \coordinate (P_S) at (1.6, 4.0);
        \draw (P_A) -- (P_B) -- (P_C) -- (P_S) -- cycle;
        \draw (P_B) -- (P_S);
        \draw[dashed] (P_A) -- (P_D) -- (P_C);
        \draw[dashed] (P_D) -- (P_S);
    \end{scope}
\end{tikzpicture}
\end{minipage}
\noindent
\begin{minipage}[c]{0.58\textwidth}
\zadnum На рисунку зображено правильну чотирикутну призму та правильну чотирикутну піраміду. Сторона основи піраміди дорівнює стороні основи призми. Усі ребра піраміди рівні й дорівнюють $2\sqrt{2}$~см. Знайдіть об'єм (у см$^3$) призми, якщо сума довжин висот призми та піраміди дорівнює $10$~см.
\par\vspace{0.4cm}
\nmtAnswerBox
\end{minipage}%
\hfill
\begin{minipage}[c]{0.40\textwidth}
\centering
\begin{tikzpicture}[scale=0.65, thick, line join=round]
    \coordinate (A) at (0,0);
    \coordinate (B) at (1.8,0);
    \coordinate (C) at (2.6, 0.7);
    \coordinate (D) at (0.8, 0.7);
    \coordinate (A1) at (0,3.8);
    \coordinate (B1) at (1.8,3.8);
    \coordinate (C1) at (2.6, 4.5);
    \coordinate (D1) at (0.8, 4.5);
    \draw (A) -- (B) -- (C) -- (C1) -- (D1) -- (A1) -- cycle;
    \draw (B) -- (B1) -- (C1);
    \draw (A1) -- (B1);
    \draw[dashed] (A) -- (D) -- (C);
    \draw[dashed] (D) -- (D1);
    \begin{scope}[shift={(3.8,0)}]
        \coordinate (P_A) at (0,0);
        \coordinate (P_B) at (2.2,0);
        \coordinate (P_C) at (3.2, 0.9);
        \coordinate (P_D) at (1.0, 0.9);
        \coordinate (P_S) at (1.6, 4.0);
        \draw (P_A) -- (P_B) -- (P_C) -- (P_S) -- cycle;
        \draw (P_B) -- (P_S);
        \draw[dashed] (P_A) -- (P_D) -- (P_C);
        \draw[dashed] (P_D) -- (P_S);
    \end{scope}
\end{tikzpicture}
\end{minipage}
\taskBlock{Призма з ромбом в основі та конус --- сесія 5.06, завдання №21}
\noindent
\begin{minipage}[c]{0.62\textwidth}
\zadnum На рисунку зображено пряму чотирикутну призму та конус. В основі призми лежить ромб, площа якого дорівнює $216$ см$^2$, а сторона~--- $15$ см. Радіус основи конуса дорівнює половині \textit{висоти} ромба. Висота призми дорівнює твірній конуса. Визначте площу \textit{повної} поверхні призми (у см$^2$), якщо площа бічної поверхні конуса дорівнює $144\pi$ см$^2$.
\end{minipage}%
\hfill
\begin{minipage}[c]{0.36\textwidth}
\centering
\begin{tikzpicture}[scale=0.7, thick, line join=round]
    \coordinate (A) at (0,0);
    \coordinate (B) at (1.8,0);
    \coordinate (C) at (2.6,0.7);
    \coordinate (D) at (0.8,0.7);
    \coordinate (A1) at (0,3.4);
    \coordinate (B1) at (1.8,3.4);
    \coordinate (C1) at (2.6,4.1);
    \coordinate (D1) at (0.8,4.1);
    \draw (A) -- (B) -- (C) -- (C1) -- (D1) -- (A1) -- cycle;
    \draw (B) -- (B1) -- (C1);
    \draw (A1) -- (B1);
    \draw[dashed] (A) -- (D) -- (C);
    \draw[dashed] (D) -- (D1);
    \begin{scope}[shift={(4.0,0)}]
        \draw[thick, dashed] (-0.7,0) arc (180:0:0.7 and 0.2);
        \draw[thick] (-0.7,0) arc (180:360:0.7 and 0.2);
        \draw[thick] (-0.7,0) -- (0,2.6) -- (0.7,0);
    \end{scope}
\end{tikzpicture}
\end{minipage}
\nmtAnswerBox
\solbtn{64}
\noindent
\begin{minipage}[c]{0.62\textwidth}
\zadnum На рисунку зображено пряму чотирикутну призму та конус. В основі призми лежить ромб зі стороною $10$ см і гострим кутом $30^\circ$. Радіус основи конуса дорівнює половині \textit{висоти} ромба. Висота призми дорівнює твірній конуса. Визначте площу \textit{повної} поверхні призми (у см$^2$), якщо площа бічної поверхні конуса дорівнює $25\pi$ см$^2$.
\end{minipage}%
\hfill
\begin{minipage}[c]{0.36\textwidth}
\centering
\begin{tikzpicture}[scale=0.7, thick, line join=round]
    \coordinate (A) at (0,0);
    \coordinate (B) at (1.4,0.5);
    \coordinate (C) at (2.6,0);
    \coordinate (D) at (1.2,-0.5);
    \draw[thick] (A) -- (B) -- (C) -- (D) -- cycle;
    \draw[thick] (A) -- (0,2.6);
    \draw[thick] (B) -- (1.4,3.1);
    \draw[thick] (C) -- (2.6,2.6);
    \draw[thick, dashed] (D) -- (1.2,2.1);
    \draw[thick] (0,2.6) -- (1.4,3.1) -- (2.6,2.6) -- (1.2,2.1) -- cycle;
    \begin{scope}[shift={(4.0,0)}]
        \draw[thick, dashed] (-0.7,0) arc (180:0:0.7 and 0.2);
        \draw[thick] (-0.7,0) arc (180:360:0.7 and 0.2);
        \draw[thick] (-0.7,0) -- (0,2.6) -- (0.7,0);
    \end{scope}
\end{tikzpicture}
\end{minipage}
\nmtAnswerBox
\noindent
\begin{minipage}[c]{0.62\textwidth}
\zadnum На рисунку зображено пряму чотирикутну призму та конус. В основі призми лежить ромб, площа якого дорівнює $120$ см$^2$, а сторона~--- $10$ см. Радіус основи конуса дорівнює половині \textit{висоти} ромба. Висота призми дорівнює твірній конуса. Визначте \textit{об'єм} призми (у см$^3$), якщо площа бічної поверхні конуса дорівнює $90\pi$ см$^2$.
\end{minipage}%
\hfill
\begin{minipage}[c]{0.36\textwidth}
\centering
\begin{tikzpicture}[scale=0.7, thick, line join=round]
    \coordinate (A) at (0,0);
    \coordinate (B) at (1.4,0.5);
    \coordinate (C) at (2.6,0);
    \coordinate (D) at (1.2,-0.5);
    \draw[thick] (A) -- (B) -- (C) -- (D) -- cycle;
    \draw[thick] (A) -- (0,2.6);
    \draw[thick] (B) -- (1.4,3.1);
    \draw[thick] (C) -- (2.6,2.6);
    \draw[thick, dashed] (D) -- (1.2,2.1);
    \draw[thick] (0,2.6) -- (1.4,3.1) -- (2.6,2.6) -- (1.2,2.1) -- cycle;
    \begin{scope}[shift={(4.0,0)}]
        \draw[thick, dashed] (-0.7,0) arc (180:0:0.7 and 0.2);
        \draw[thick] (-0.7,0) arc (180:360:0.7 and 0.2);
        \draw[thick] (-0.7,0) -- (0,2.6) -- (0.7,0);
    \end{scope}
\end{tikzpicture}
\end{minipage}
\nmtAnswerBox
\chapterTitle{РОЗВ\textquotesingle ЯЗКИ ОРИГІНАЛЬНИХ ЗАДАЧ}
\noindent\small У цьому розділі зібрано стислі розв\textquotesingle язки до перших задач кожного блоку (оригіналів НМТ-2026). Натискайте посилання \fbox{\footnotesize $\leftarrow$ Задача №$N$}, щоб повернутися до задачі.
\par\vspace{0.4cm}
\par\vspace{0.3cm}\noindent\hypertarget{sol:1}{}
\begin{tcolorbox}[colback=titleBg!40, colframe=mainGreen!55, boxrule=0.4pt, arc=2pt, left=6pt, right=6pt, top=4pt, bottom=4pt, breakable]
\noindent\small\textbf{Задача №1.}\quad Спільними точками площин $ABS$ і $MCS$ є $S$ (вершина) та $M$ (бо $M\in AB\subset ABS$ і $M\in MCS$). Через дві спільні точки площини перетинаються по прямій $SM$. Відповідь: \textbf{Д}.
\par\vspace{0.1cm}\hfill\hyperlink{task:1}{\fbox{\footnotesize $\leftarrow$ Задача №1}}
\end{tcolorbox}
\par\vspace{0.3cm}\noindent\hypertarget{sol:4}{}
\begin{tcolorbox}[colback=titleBg!40, colframe=mainGreen!55, boxrule=0.4pt, arc=2pt, left=6pt, right=6pt, top=4pt, bottom=4pt, breakable]
\noindent\small\textbf{Задача №4.}\quad Апофема~--- це висота бічної грані, проведена з вершини $S$. Правильна трикутна піраміда має $3$ бічні грані, отже $3$ апофеми. Відповідь: \textbf{Г} (лише три).
\par\vspace{0.1cm}\hfill\hyperlink{task:4}{\fbox{\footnotesize $\leftarrow$ Задача №4}}
\end{tcolorbox}
\par\vspace{0.3cm}\noindent\hypertarget{sol:7}{}
\begin{tcolorbox}[colback=titleBg!40, colframe=mainGreen!55, boxrule=0.4pt, arc=2pt, left=6pt, right=6pt, top=4pt, bottom=4pt, breakable]
\noindent\small\textbf{Задача №7.}\quad Призма має $12$ ребер. Із прямою $B_1D$ перетинаються по $3$ ребра біля кожної з вершин $B_1$ і $D$ (разом $6$); решта $6$ ребер не мають із $B_1D$ спільних точок і не паралельні їй, тобто мимобіжні. Відповідь: \textbf{В}.
\par\vspace{0.1cm}\hfill\hyperlink{task:7}{\fbox{\footnotesize $\leftarrow$ Задача №7}}
\end{tcolorbox}
\par\vspace{0.3cm}\noindent\hypertarget{sol:10}{}
\begin{tcolorbox}[colback=titleBg!40, colframe=mainGreen!55, boxrule=0.4pt, arc=2pt, left=6pt, right=6pt, top=4pt, bottom=4pt, breakable]
\noindent\small\textbf{Задача №10.}\quad $AB_1$ --- діагональ бічної грані. Паралельна їй пряма через дві вершини куба лише одна --- це діагональ протилежної грані $DC_1$. Відповідь: \textbf{Б}.
\par\vspace{0.1cm}\hfill\hyperlink{task:10}{\fbox{\footnotesize $\leftarrow$ Задача №10}}
\end{tcolorbox}
\par\vspace{0.3cm}\noindent\hypertarget{sol:13}{}
\begin{tcolorbox}[colback=titleBg!40, colframe=mainGreen!55, boxrule=0.4pt, arc=2pt, left=6pt, right=6pt, top=4pt, bottom=4pt, breakable]
\noindent\small\textbf{Задача №13.}\quad Конус дістаємо обертанням прямокутного трикутника навколо одного з його катетів: цей катет стає віссю (висотою), другий катет~--- радіусом основи, гіпотенуза~--- твірною. Відповідь: \textbf{Г}.
\par\vspace{0.1cm}\hfill\hyperlink{task:13}{\fbox{\footnotesize $\leftarrow$ Задача №13}}
\end{tcolorbox}
\par\vspace{0.3cm}\noindent\hypertarget{sol:16}{}
\begin{tcolorbox}[colback=titleBg!40, colframe=mainGreen!55, boxrule=0.4pt, arc=2pt, left=6pt, right=6pt, top=4pt, bottom=4pt, breakable]
\noindent\small\textbf{Задача №16.}\quad Призма, конус, циліндр і піраміда мають основу (основи). Куля основи не має. Відповідь: \textbf{Г}.
\par\vspace{0.1cm}\hfill\hyperlink{task:16}{\fbox{\footnotesize $\leftarrow$ Задача №16}}
\end{tcolorbox}
\par\vspace{0.3cm}\noindent\hypertarget{sol:19}{}
\begin{tcolorbox}[colback=titleBg!40, colframe=mainGreen!55, boxrule=0.4pt, arc=2pt, left=6pt, right=6pt, top=4pt, bottom=4pt, breakable]
\noindent\small\textbf{Задача №19.}\quad Площина має зі сферою лише одну спільну точку тоді й тільки тоді, коли вона дотична до сфери. Через точку дотику $M$ можна провести безліч площин, перпендикулярних до $\beta$, і кожна з них дотикається сфери в $M$. Отже, таких площин безліч. Відповідь: \textbf{Д}.
\par\vspace{0.1cm}\hfill\hyperlink{task:19}{\fbox{\footnotesize $\leftarrow$ Задача №19}}
\end{tcolorbox}
\par\vspace{0.3cm}\noindent\hypertarget{sol:22}{}
\begin{tcolorbox}[colback=titleBg!40, colframe=mainGreen!55, boxrule=0.4pt, arc=2pt, left=6pt, right=6pt, top=4pt, bottom=4pt, breakable]
\noindent\small\textbf{Задача №22.}\quad $|\vec{a}|=\sqrt{5^2+2^2+(-1)^2}=\sqrt{25+4+1}=\sqrt{30}$. Відповідь: \textbf{В}.
\par\vspace{0.1cm}\hfill\hyperlink{task:22}{\fbox{\footnotesize $\leftarrow$ Задача №22}}
\end{tcolorbox}
\par\vspace{0.3cm}\noindent\hypertarget{sol:25}{}
\begin{tcolorbox}[colback=titleBg!40, colframe=mainGreen!55, boxrule=0.4pt, arc=2pt, left=6pt, right=6pt, top=4pt, bottom=4pt, breakable]
\noindent\small\textbf{Задача №25.}\quad Найбільша відстань між точками сфери~--- це діаметр, тож $MN$ є діаметром. $MN=\sqrt{(4-(-4))^2+(-3-5)^2+(4-0)^2}=\sqrt{64+64+16}=\sqrt{144}=12$. Радіус $R=\dfrac{12}{2}=6$. Відповідь: \textbf{Б} ($6$).
\par\vspace{0.1cm}\hfill\hyperlink{task:25}{\fbox{\footnotesize $\leftarrow$ Задача №25}}
\end{tcolorbox}
\par\vspace{0.3cm}\noindent\hypertarget{sol:28}{}
\begin{tcolorbox}[colback=titleBg!40, colframe=mainGreen!55, boxrule=0.4pt, arc=2pt, left=6pt, right=6pt, top=4pt, bottom=4pt, breakable]
\noindent\small\textbf{Задача №28.}\quad Відстань від $O$ до $(x;y;z)$ дорівнює $\sqrt{x^2+y^2+z^2}$. Для чотирьох точок маємо $\sqrt{100}=10$, а для $(9;1;0)$: $\sqrt{81+1}=\sqrt{82}\approx 9{,}06\neq 10$. Відповідь: \textbf{В}.
\par\vspace{0.1cm}\hfill\hyperlink{task:28}{\fbox{\footnotesize $\leftarrow$ Задача №28}}
\end{tcolorbox}
\par\vspace{0.3cm}\noindent\hypertarget{sol:31}{}
\begin{tcolorbox}[colback=titleBg!40, colframe=mainGreen!55, boxrule=0.4pt, arc=2pt, left=6pt, right=6pt, top=4pt, bottom=4pt, breakable]
\noindent\small\textbf{Задача №31.}\quad Оскільки $\vec{a}=\overrightarrow{KM}=M-K$, маємо $K=M-\vec{a}=(-7-5;\,1-(-3);\,4-8)=(-12;\,4;\,-4)$. Відповідь: \textbf{Б}.
\par\vspace{0.1cm}\hfill\hyperlink{task:31}{\fbox{\footnotesize $\leftarrow$ Задача №31}}
\end{tcolorbox}
\par\vspace{0.3cm}\noindent\hypertarget{sol:34}{}
\begin{tcolorbox}[colback=titleBg!40, colframe=mainGreen!55, boxrule=0.4pt, arc=2pt, left=6pt, right=6pt, top=4pt, bottom=4pt, breakable]
\noindent\small\textbf{Задача №34.}\quad $\vec{c} = \vec{a} + \vec{b} = (-2+4;\, 2-3;\, 1+0) = (2;\, -1;\, 1)$. Тоді $|\vec{c}| = \sqrt{2^2 + (-1)^2 + 1^2} = \sqrt{6}$. Відповідь: \textbf{Г} ($\sqrt6$).
\par\vspace{0.1cm}\hfill\hyperlink{task:34}{\fbox{\footnotesize $\leftarrow$ Задача №34}}
\end{tcolorbox}
\par\vspace{0.3cm}\noindent\hypertarget{sol:37}{}
\begin{tcolorbox}[colback=titleBg!40, colframe=mainGreen!55, boxrule=0.4pt, arc=2pt, left=6pt, right=6pt, top=4pt, bottom=4pt, breakable]
\noindent\small\textbf{Задача №37.}\quad Твірна дорівнює відстані між центрами основ: $O_1O_2=\sqrt{(3+3)^2+(4-2)^2+(4-1)^2}=\sqrt{36+4+9}=\sqrt{49}=7$. Відповідь: \textbf{А}.
\par\vspace{0.1cm}\hfill\hyperlink{task:37}{\fbox{\footnotesize $\leftarrow$ Задача №37}}
\end{tcolorbox}
\par\vspace{0.3cm}\noindent\hypertarget{sol:40}{}
\begin{tcolorbox}[colback=titleBg!40, colframe=mainGreen!55, boxrule=0.4pt, arc=2pt, left=6pt, right=6pt, top=4pt, bottom=4pt, breakable]
\noindent\small\textbf{Задача №40.}\quad $\overrightarrow{AB}=(3-1;\,1-2;\,1-3)=(2;\,-1;\,-2)$. $|\overrightarrow{AB}|=\sqrt{2^2+(-1)^2+(-2)^2}=\sqrt{9}=3$. Відповідь: \textbf{А}.
\par\vspace{0.1cm}\hfill\hyperlink{task:40}{\fbox{\footnotesize $\leftarrow$ Задача №40}}
\end{tcolorbox}
\par\vspace{0.3cm}\noindent\hypertarget{sol:43}{}
\begin{tcolorbox}[colback=titleBg!40, colframe=mainGreen!55, boxrule=0.4pt, arc=2pt, left=6pt, right=6pt, top=4pt, bottom=4pt, breakable]
\noindent\small\textbf{Задача №43.}\quad Координати середини~--- середні арифметичні: $\left(\dfrac{1+3}{2};\,\dfrac{2+1}{2};\,\dfrac{3+1}{2}\right)=(2;\,1{,}5;\,2)$. Відповідь: \textbf{В}.
\par\vspace{0.1cm}\hfill\hyperlink{task:43}{\fbox{\footnotesize $\leftarrow$ Задача №43}}
\end{tcolorbox}
\par\vspace{0.3cm}\noindent\hypertarget{sol:46}{}
\begin{tcolorbox}[colback=titleBg!40, colframe=mainGreen!55, boxrule=0.4pt, arc=2pt, left=6pt, right=6pt, top=4pt, bottom=4pt, breakable]
\noindent\small\textbf{Задача №46.}\quad Радіус конуса $r=\sqrt{17^2-15^2}=\sqrt{64}=8$ см, тож радіус циліндра $R=2r=16$ см. З бічної поверхні $2\pi R h=640\pi$ маємо $h=\dfrac{640}{2\cdot 16}=20$ см. Об'єм $V=\pi R^2 h=\pi\cdot 256\cdot 20=5120\pi$, тож $\dfrac{V}{\pi}=5120$. Відповідь: \textbf{5120}.
\par\vspace{0.1cm}\hfill\hyperlink{task:46}{\fbox{\footnotesize $\leftarrow$ Задача №46}}
\end{tcolorbox}
\par\vspace{0.3cm}\noindent\hypertarget{sol:49}{}
\begin{tcolorbox}[colback=titleBg!40, colframe=mainGreen!55, boxrule=0.4pt, arc=2pt, left=6pt, right=6pt, top=4pt, bottom=4pt, breakable]
\noindent\small\textbf{Задача №49.}\quad У трапецію з вписаним колом сума основ дорівнює сумі бічних сторін: $a+b=2\cdot26=52$. Площа основи $\dfrac{a+b}{2}\cdot h_{\text{тр}}=624\Rightarrow 26\,h_{\text{тр}}=624\Rightarrow h_{\text{тр}}=24$~--- це діаметр основи конуса, тож $r=12$. Висота конуса (= висота призми) $h=\sqrt{20^2-12^2}=\sqrt{256}=16$. Периметр трапеції $P=(a+b)+2\cdot26=52+52=104$. Бічна поверхня призми $P\cdot h=104\cdot16=1664$. Відповідь: $1664$.
\par\vspace{0.1cm}\hfill\hyperlink{task:49}{\fbox{\footnotesize $\leftarrow$ Задача №49}}
\end{tcolorbox}
\par\vspace{0.3cm}\noindent\hypertarget{sol:52}{}
\begin{tcolorbox}[colback=titleBg!40, colframe=mainGreen!55, boxrule=0.4pt, arc=2pt, left=6pt, right=6pt, top=4pt, bottom=4pt, breakable]
\noindent\small\textbf{Задача №52.}\quad Для основи піраміди $r=\dfrac{a}{2\sqrt3}=4\sqrt3$, звідки сторона $a=24$. Площа бічної поверхні $S=\dfrac12\cdot 3a\cdot l=288$, тож $l=\dfrac{288}{36}=8$ (апофема). З $l^2=h^2+r^2$: $h^2=64-48=16$, $h=4$. Тоді висота й радіус конуса: $h=4$, $R=4\sqrt3$, а $V=\dfrac13\pi R^2 h=\dfrac13\pi\cdot48\cdot4=64\pi$. Отже $\dfrac{V}{\pi}=64$. Відповідь: \textbf{$64$}.
\par\vspace{0.1cm}\hfill\hyperlink{task:52}{\fbox{\footnotesize $\leftarrow$ Задача №52}}
\end{tcolorbox}
\par\vspace{0.3cm}\noindent\hypertarget{sol:55}{}
\begin{tcolorbox}[colback=titleBg!40, colframe=mainGreen!55, boxrule=0.4pt, arc=2pt, left=6pt, right=6pt, top=4pt, bottom=4pt, breakable]
\noindent\small\textbf{Задача №55.}\quad Конус: $r:l=3:5$, $S_{\text{бік}}=\pi rl=135\pi$. Із $l=\dfrac{5r}{3}$: $\pi r\cdot\dfrac{5r}{3}=135\pi$, тож $r^2=81$, $r=9$, $l=15$, висота $h=\sqrt{l^2-r^2}=\sqrt{225-81}=12$, $V_{\text{кон}}=\dfrac13\pi r^2h=\dfrac13\pi\cdot81\cdot12=324\pi$. Півсфера: $3\pi R^2=243\pi$, $R=9$, $V_{\text{півсф}}=\dfrac23\pi R^3=\dfrac23\pi\cdot729=486\pi$. Відношення $\dfrac{486\pi}{324\pi}=1{,}5$.
\par\vspace{0.1cm}\hfill\hyperlink{task:55}{\fbox{\footnotesize $\leftarrow$ Задача №55}}
\end{tcolorbox}
\par\vspace{0.3cm}\noindent\hypertarget{sol:58}{}
\begin{tcolorbox}[colback=titleBg!40, colframe=mainGreen!55, boxrule=0.4pt, arc=2pt, left=6pt, right=6pt, top=4pt, bottom=4pt, breakable]
\noindent\small\textbf{Задача №58.}\quad Діагональ основи паралелепіпеда $d = \sqrt{9^2+12^2} = \sqrt{225} = 15$ см~--- це діаметр циліндра, тож $r = 7{,}5$ см. Висота паралелепіпеда $H = \sqrt{39^2 - 15^2} = \sqrt{1296} = 36$ см, а висота циліндра $h = \dfrac23 H = 24$ см. Об'єм $V = \pi r^2 h = \pi\cdot 7{,}5^2\cdot 24 = 1350\pi$, тож $\dfrac{V}{\pi} = 1350$. Відповідь: \textbf{$1350$}.
\par\vspace{0.1cm}\hfill\hyperlink{task:58}{\fbox{\footnotesize $\leftarrow$ Задача №58}}
\end{tcolorbox}
\par\vspace{0.3cm}\noindent\hypertarget{sol:61}{}
\begin{tcolorbox}[colback=titleBg!40, colframe=mainGreen!55, boxrule=0.4pt, arc=2pt, left=6pt, right=6pt, top=4pt, bottom=4pt, breakable]
\noindent\small\textbf{Задача №61.}\quad Усі ребра піраміди $=3\sqrt2$, тож сторона її основи $a_p=3\sqrt2$, а діагональ основи $a_p\sqrt2=6$. Висота піраміди $h_2=\sqrt{(3\sqrt2)^2-3^2}=\sqrt{18-9}=3$. Сторона основи піраміди вдвічі менша від діагоналі основи призми, тож діагональ призми $=6\sqrt2$, а її сторона $a=\dfrac{6\sqrt2}{\sqrt2}=6$, площа основи $S=36$. Висота призми $h_1=13-3=10$, об'єм $V=36\cdot10=360$~см$^3$. Відповідь: \textbf{360}.
\par\vspace{0.1cm}\hfill\hyperlink{task:61}{\fbox{\footnotesize $\leftarrow$ Задача №61}}
\end{tcolorbox}
\par\vspace{0.3cm}\noindent\hypertarget{sol:64}{}
\begin{tcolorbox}[colback=titleBg!40, colframe=mainGreen!55, boxrule=0.4pt, arc=2pt, left=6pt, right=6pt, top=4pt, bottom=4pt, breakable]
\noindent\small\textbf{Задача №64.}\quad Висота ромба: $S=a\cdot h\Rightarrow h=\dfrac{216}{15}=14{,}4$. Радіус конуса $r=\dfrac{h}{2}=7{,}2$. З $\pi r l=144\pi$ маємо $l=\dfrac{144}{7{,}2}=20$~--- це й висота призми $H$. Повна поверхня призми $=2S+P\cdot H=2\cdot216+(4\cdot15)\cdot20=432+1200=1632$ см$^2$. Відповідь: \textbf{$1632$}.
\par\vspace{0.1cm}\hfill\hyperlink{task:64}{\fbox{\footnotesize $\leftarrow$ Задача №64}}
\end{tcolorbox}
\chapterTitle{ВІДПОВІДІ}
\noindent\small Відповіді наведено у тому ж порядку, що й завдання.\par\vspace{0.3cm}
\ansType{Завдання з вибором однієї відповіді (А--Д)}
\ansTheme{Просторові фігури та їх властивості}
\begin{multicols}{3}\noindent
\ansitem{Д}\ansitem{В}\ansitem{Г}
\ansitem{Г}\ansitem{Г}\ansitem{Д}
\ansitem{В}\ansitem{Б}\ansitem{В}
\ansitem{Б}\ansitem{А}\ansitem{Г}
\ansitem{Г}\ansitem{В}\ansitem{Б}
\ansitem{Г}\ansitem{В}\ansitem{Д}
\ansitem{Д}\ansitem{Б}\ansitem{Д}
\end{multicols}
\ansTheme{Вектори і координати в просторі}
\begin{multicols}{3}\noindent
\ansitem{В}\ansitem{Г}\ansitem{В}
\ansitem{Б}\ansitem{Б}\ansitem{В}
\ansitem{В}\ansitem{Б}\ansitem{В}
\ansitem{Г}\ansitem{В}\ansitem{Б}
\ansitem{Г}\ansitem{Б}\ansitem{А}
\ansitem{А}\ansitem{Г}\ansitem{А}
\ansitem{А}\ansitem{Б}\ansitem{Б}
\ansitem{В}\ansitem{Г}\ansitem{А}
\end{multicols}
\ansType{Завдання з короткою відповіддю}
\begin{multicols}{3}\noindent
\ansitem{$5120$}\ansitem{$6912$}\ansitem{$36$}
\ansitem{$1664$}\ansitem{$576$}\ansitem{$240$}
\ansitem{$64$}\ansitem{$24$}\ansitem{$108$}
\ansitem{$1{,}5$}\ansitem{$2$}\ansitem{$6$}
\ansitem{$1350$}\ansitem{$250$}\ansitem{$2500$}
\ansitem{$360$}\ansitem{$720$}\ansitem{$64$}
\ansitem{$1632$}\ansitem{$500$}\ansitem{$1800$}
\end{multicols}
\end{document}
